\documentclass[10pt]{article}

% ---------- Document Identity (update these only) ----------
\newcommand{\DocStage}{}
\newcommand{\DocTitle}{Your Road to Tier One SOF}
\newcommand{\ProgramName}{182-Week Civilian-to-Selection Readiness Pipeline}
\newcommand{\ProgramSubtitle}{}

% ---------- Page ----------
\usepackage[legalpaper,left=0.5in,right=0.5in,top=0.6in,bottom=0.45in]{geometry}

% ---------- Typography ----------
\usepackage[T1]{fontenc}
\usepackage[utf8]{inputenc}
\usepackage{libertinus}
\usepackage{microtype}
\usepackage{parskip}
\usepackage{ragged2e}
\usepackage{textcomp}
\AtBeginDocument{\color{black}}
\AtBeginDocument{\pagecolor{white}}
\AtBeginDocument{\justifying}

% ---------- Landscape pages ----------
\usepackage{pdflscape}

% ---------- Color ----------
\usepackage[table]{xcolor}
\definecolor{StageBlue}{HTML}{1F4E79}
\definecolor{StageBlueDark}{HTML}{163A5A}
\definecolor{LightGray}{HTML}{F2F4F7}
\definecolor{MidGray}{HTML}{D9DEE5}
\definecolor{RuleGray}{HTML}{C7CED8}
\definecolor{HeaderInk}{HTML}{000000}
\definecolor{Ink}{HTML}{000000}

% ---------- Utility ----------
\usepackage{enumitem}
\sloppy
\setlength{\emergencystretch}{2em}

% ---------- Boxes / UI ----------
\usepackage[most]{tcolorbox}

\tcbset{
  charterTitle/.style={
    enhanced,
    colback=StageBlue!4,
    colframe=StageBlue,
    boxrule=1.25pt,
    arc=3pt,
    left=16pt,right=16pt,top=14pt,bottom=14pt,
    borderline north={3.2pt}{0pt}{StageBlue},
    borderline west={4pt}{0pt}{StageBlue},
    borderline south={1.25pt}{0pt}{StageBlue},
    drop shadow={black!25!white},
  },
  charterRule/.style={
    enhanced,
    colback=LightGray,
    colframe=RuleGray,
    boxrule=0.85pt,
    arc=3pt,
    left=12pt,right=12pt,top=10pt,bottom=10pt,
    borderline west={3pt}{0pt}{StageBlue},
  },
  charterBadge/.style={
    enhanced,
    colback=StageBlue,
    coltext=white,
    colframe=StageBlueDark,
    boxrule=0.6pt,
    arc=2pt,
    left=6pt,right=6pt,top=3pt,bottom=3pt,
  }
}
\newtcbox{\Badge}{nobeforeafter,charterBadge}

% ---------- Lists ----------
\setlist[itemize]{leftmargin=1.5em, itemsep=0.25em, topsep=0.25em}
\setlist[enumerate]{leftmargin=1.8em, itemsep=0.25em, topsep=0.25em}

% ---------- TOC ----------
\usepackage{tocloft}
\setcounter{tocdepth}{3}
\setlength{\cftbeforesecskip}{3pt}
\setlength{\cftbeforesubsecskip}{2pt}
\setlength{\cftbeforesubsubsecskip}{0pt}
\renewcommand{\cftsecfont}{\bfseries}
\renewcommand{\cftsecpagefont}{\bfseries}
\renewcommand{\cftsubsecfont}{\normalfont}
\renewcommand{\cftsubsecpagefont}{\normalfont}
\renewcommand{\cftsubsubsecfont}{\normalfont}
\renewcommand{\cftsubsubsecpagefont}{\normalfont}
\setlength{\cftsecindent}{0pt}
\setlength{\cftsubsecindent}{1.25em}
\setlength{\cftsubsubsecindent}{2.8em}
\setlength{\cftsecnumwidth}{2.1em}
\setlength{\cftsubsecnumwidth}{2.8em}
\setlength{\cftsubsubsecnumwidth}{3.6em}

% Remove the built-in TOC heading so only the custom heading is shown
\makeatletter
\renewcommand{\tableofcontents}{\@starttoc{toc}}
\makeatother

% ---------- Header/Footer: page number only ----------
\usepackage{fancyhdr}
\pagestyle{fancy}
\fancyhf{}
\renewcommand{\headrulewidth}{0pt}
\renewcommand{\footrulewidth}{0pt}
\fancyfoot[C]{\thepage}

% ---------- Section formatting ----------
\usepackage{titlesec}

\titleformat{\section}
  {\Large\bfseries\color{Ink}}
  {\thesection}{0.6em}{}
  [\vspace{0.25em}\textcolor{StageBlue}{\rule{\linewidth}{0.8pt}}]

\titleformat{\subsection}
  {\large\bfseries\color{Ink}}
  {\thesubsection}{0.6em}{}

\titleformat{\subsubsection}
  {\normalsize\bfseries\color{Ink}}
  {\thesubsubsection}{0.6em}{}

\titlespacing*{\section}{0pt}{0.95em}{0.35em}
\titlespacing*{\subsection}{0pt}{0.65em}{0.25em}
\titlespacing*{\subsubsection}{0pt}{0.55em}{0.2em}

% ---------- Table engine ----------
\usepackage{tabularray}
\UseTblrLibrary{booktabs}

% ---------- tabularray: GLOBAL table treatment ----------
\SetTblrInner{
  cells = {fg=black, bg=white, valign=t},
}

% ---------- Header helpers ----------
\newcommand{\Hdr}[1]{\shortstack[c]{\strut #1\strut}}

% ---------- Lines ----------
\newcommand{\TitleLine}{\textcolor{StageBlue}{\rule{0.98\linewidth}{0.95pt}}}
\newcommand{\SubTitleLine}{\textcolor{RuleGray}{\rule{0.98\linewidth}{0.65pt}}}

% ---------- Continued on next page ----------
\newcommand{\ContinuedOnNextPageInline}{%
  \par
  \vfill
  \noindent\textcolor{RuleGray}{\rule{\linewidth}{1pt}}\vspace{-2pt}
  \noindent\hfill\textit{Continued on next page}\par\vspace{-10pt}
  \noindent\textcolor{RuleGray}{\rule{\linewidth}{1pt}}
  \clearpage
}

% ---------- PDF hyperlinks ----------
\usepackage[hidelinks]{hyperref}
\hypersetup{
  colorlinks=false,
  pdfborder={0 0 0},
  linktoc=all,
  pdftitle={\DocTitle},
  pdfauthor={\ProgramName}
}

% =========================================================
\begin{document}

\section*{7.10 — Logistics, Maintenance, Storage, and Sustainment Equipment}
\phantomsection
\addcontentsline{toc}{section}{7.10 — Logistics, Maintenance, Storage, and Sustainment Equipment}

Overview \textbf{ID}: 7.10\\
Overview Name: Logistics, Maintenance, Storage, and Sustainment Equipment\\
Version: v1.0\\
Governing Status: Draft — Non-Deployable\\
Scope Boundary Statement: Covers storage systems, organization controls, drying workflows, maintenance tooling, and sustainment posture required to keep equipment usable and safe; does not provide programming, protocols, calendars, medical guidance, or brand / retailer links.

\subsection*{7.10.1 Purpose \& Scope}
\phantomsection
\addcontentsline{toc}{subsection}{7.10.1 Purpose \& Scope}

7.10 covers storage systems, organization posture, and maintenance tooling that keep training equipment usable and safe, including drying and sustainment workflows and staging checklists that prevent missed sessions and unsafe improvisation.

This overview crosswalks Stage 6 category timing to the logistics and sustainment conditions required for safe execution and admissible evidence.

It does not cover, except by cross-reference:

\begin{itemize}
\item training tools: Stage 7.01, Stage 7.07, Stage 7.08
\item measurement / timing tools: Stage 7.18
\item hygiene / health consumables: Stage 7.20
\item apparel itself: Stage 7.14, Stage 7.15
\end{itemize}

This overview does not provide programming, test protocols, calendars, medical guidance, or brand / retailer links.

\subsection*{7.10.2 Governance And Safety Boundaries}
\phantomsection
\addcontentsline{toc}{subsection}{7.10.2 Governance And Safety Boundaries}

\begin{itemize}
\item Stage 0 boundary alignment: stop rules, safety override doctrine, conflict-resolution hierarchy, and evidence-admissibility controls apply to all 7.10 selections and substitutions.
\item Sustainment enables adherence: logistics prevents missed sessions and unsafe improvisation.
\item Safety priorities:
  \begin{itemize}
  \item eliminate trip and tip hazards,
  \item prevent equipment failure via inspection and timely repair/replacement,
  \item keep gear dry to reduce mold and skin irritation risk.
  \end{itemize}
\item Start-from-zero constraints:
  \begin{itemize}
  \item implement low-cost organization early and scale only as inventory grows.
  \end{itemize}
\item Validity posture:
  \begin{itemize}
  \item consistent setup and tool readiness supports comparable execution and test-day readiness.
  \end{itemize}
\item Escalation posture: hold use, substitute, and correct immediately when shelving is unstable, drying control fails, mold is present, sharp damage is observed, or clutter creates a trip path.
\item Evidence inadmissibility: evidence produced under unsafe, unstable, wet, damaged, or improvised sustainment conditions is inadmissible for readiness claims until corrected and re-verified.
\item Legal/ethical sourcing only; no illegal procurement guidance.
\item No medical advice; direct medical items to Stage 7.20 by cross-reference only.
\end{itemize}

\subsection*{7.10.3 Tier Doctrine}
\phantomsection
\addcontentsline{toc}{subsection}{7.10.3 Tier Doctrine}

\begin{itemize}
\item Price band convention: \$ = minimal household-cost item; \$\$ = modest purpose-built item; \$\$\$ = expanded-capability or redundancy item.
\item N/A rule: use N/A only where a price band genuinely does not apply and omission does not obscure safety, maintenance burden, or phase tie-in logic.
\item Substitution posture: substitutions are acceptable only when they preserve storage stability, drying control, safe access, and maintenance intent; substitutions that increase clutter, moisture retention, instability, or failure risk break intent.
\item Tier 0: household-only storage and ad hoc organization; minimal maintenance capacity; sufficient only for minimal inventory and low complexity.
\item Tier 1: minimum viable organization + drying + basic hand tools to safely sustain routine phase demands and prevent avoidable misses.
\item Tier 2: expanded staging and maintenance capability as inventory volume and outdoor exposure increase; includes optional go-bag/vehicle staging to prevent missed sessions.
\item Tier 3: overbuilt redundancy and expanded maintenance capability; may include light power-tool support for safe assembly/maintenance but is not required unless phase-dependent.
\end{itemize}

\subsection*{7.10.4 Selection Criteria \& Fit Checks}
\phantomsection
\addcontentsline{toc}{subsection}{7.10.4 Selection Criteria \& Fit Checks}

Fit and safety checklist (operational):

\begin{itemize}
\item Space mapping:
  \begin{itemize}
  \item define a clean zone (dry gear, measurement tools) and a wet/dirty zone (post-session gear).
  \end{itemize}
\item Stability and load rating:
  \begin{itemize}
  \item shelving must be load-rated and tip-resistant; anchor where necessary; store heavy items low.
  \end{itemize}
\item Moisture management:
  \begin{itemize}
  \item airflow and drainage; wet gear never stored sealed while damp; cross-reference Stage 7.14 / Stage 7.20 for apparel / hygiene details.
  \end{itemize}
\item Portability:
  \begin{itemize}
  \item bins and staging kits should fit vehicle and reduce “forgotten item” misses.
  \end{itemize}
\item Safety fit checks:
  \begin{itemize}
  \item no sharp edges at foot/hand height; no overloaded shelving; no blocked egress routes.
  \end{itemize}
\end{itemize}

\subsection*{7.10.5 Setup, Safety, and Anchoring}
\phantomsection
\addcontentsline{toc}{subsection}{7.10.5 Setup, Safety, and Anchoring}

Setup doctrine (no programming):

\begin{itemize}
\item Anchor tall shelving where required; maintain clear walk paths; store heavy items low.
\item Drying workflow:
  \begin{itemize}
  \item immediate post-session gear separation and drying; prevent wet footwear storage loops that create foot irritation risk.
  \end{itemize}
\item Staging workflow:
  \begin{itemize}
  \item pre-pack checklists for ruck/run/swim days (logistics only; no programming).
  \item maintain a “test week readiness” staging posture: measurement tools and required accessories staged before protected windows.
  \end{itemize}
\end{itemize}

\subsection*{7.10.6 Maintenance, Replacement, and Storage}
\phantomsection
\addcontentsline{toc}{subsection}{7.10.6 Maintenance, Replacement, and Storage}

\begin{itemize}
\item Inspection cadence:
  \begin{itemize}
  \item weekly: straps/buckles/fasteners on storage and staging systems; check trip hazards.
  \item monthly: corrosion/mold indicators, bin integrity, shelving stability.
  \end{itemize}
\item Replacement cues:
  \begin{itemize}
  \item frayed straps, cracked bins, rusted hardware, unstable shelving = remove from use and replace.
  \end{itemize}
\item Documentation:
  \begin{itemize}
  \item maintain a simple inventory list and record maintenance actions; log replacements to support auditability and reduce “mystery changes” near gates.
  \end{itemize}
\end{itemize}

\subsection*{7.10.7 Substitution Rules — Minimal-Path Alternatives}
\phantomsection
\addcontentsline{toc}{subsection}{7.10.7 Substitution Rules — Minimal-Path Alternatives}

\begin{itemize}
\item If no dedicated storage exists:
  \begin{itemize}
  \item use a single-bin essentials system plus clear zone/wet zone separation.
  \end{itemize}
\item If no drying rack/line exists:
  \begin{itemize}
  \item use a designated drying corner with airflow and strict wet gear separation.
  \end{itemize}
\item If no tool kit exists:
  \begin{itemize}
  \item treat basic hand tools as Tier 1 enabling capability and prioritize acquisition before assembling or escalating complex equipment.
  \end{itemize}
\item Stage 2.E linkage (high-level):
  \begin{itemize}
  \item missing inventory fields recorded as \textbf{MISSING}/\textbf{UNKNOWN}; no inference.
  \item tool changes and replacements logged with timestamps to preserve audit traceability.
  \item when missing logistics causes a deviation that affects admissibility (missed setup, missing required tool), treat gate evidence as compromised and reslot rather than force.
  \end{itemize}
\end{itemize}

\subsection*{7.10.8 Phase Integration Map — Required}
\phantomsection
\addcontentsline{toc}{subsection}{7.10.8 Phase Integration Map — Required}

\begingroup
\scriptsize
\setlength{\tabcolsep}{2pt}
\renewcommand{\arraystretch}{1.10}

\begin{longtblr}{
  width=\linewidth,
  rowhead=1,
  colspec = {
    Q[l,wd=1.05in]
    Q[l,wd=1.55in]
    X[l,2.75]
    X[l,3.65]
  },
  hlines,
  vlines,
  cells = {hsep=2pt, vsep=6pt, halign=l, valign=t},
  row{1} = {bg=StageBlue!15, fg=black, font=\bfseries, halign=l, valign=m},
  row{odd} = {bg=white},
  row{even} = {bg=LightGray},
}
\Hdr{Phase} &
\Hdr{Required Tier (from Stage 6)} &
\Hdr{What Changes / Becomes Mandatory} &
\Hdr{Notes} \\
Phase 00 &
Tier 1 &
Tier 1 minimal storage/drying and basic hand tools become mandatory by Phase 00 (End) &
Phase 00 (End): Tier 0 household-only examples should exist immediately. Phase 00 (End): Tier 1 sustainment reduces missed sessions and wet-gear breakdown. \\
Phase 01 &
Tier 1 &
Maintain Tier 1; organization must support increased inventory (Stage 7.01 and Stage 7.08 begin to appear) &
Phase 01 (End): avoid last-minute substitutions by staging and labeling; treat tool changes as controlled variables near gates. \\
Phase 02 &
Tier 1 &
Maintain Tier 1; assembly/maintenance becomes more consequential as strength capability grows &
Phase 02 (E/M): keep tools and fasteners available; do not assemble complex gear inside protected windows. \\
Phase 03 &
Tier 1 &
Maintain Tier 1; ruck systems introduced; spare parts posture becomes more important &
Phase 03 (End): ruck interface stability depends on sustainment and spares; avoid “broken buckle before gate” events. \\
Phase 04 &
Tier 1 &
Maintain Tier 1; aquatic category begins to appear; wet gear handling becomes more frequent &
Phase 04 (End): keep wet gear separation strict; cross-reference Stage 7.14 / Stage 7.20 for apparel / hygiene. \\
Phase 05 &
Tier 1 &
Maintain Tier 1; endurance windows increase consequence of missed sessions &
Phase 05 (End): staging discipline prevents avoidable cancellations and unsafe improvisation. \\
Phase 06 &
Tier 2 &
Tier 2 staging station and expanded spares become mandatory by Phase 06 (End) &
Phase 06 (End): expanded inventory and higher consequence windows justify Tier 2; go-bag concept reduces misses during travel/errands. \\
Phase 07 &
Tier 2 &
Maintain Tier 2; hybrid density amplifies friction costs &
Phase 07 (End): treat staging discipline as an execution control; prevent clutter and trip hazards as inventory increases. \\
Phase 08 &
Tier 2 &
Maintain Tier 2; load tolerance increases wet/dirty gear volume &
Phase 08 (End): drying and storage discipline prevents skin irritation and mold loops. \\
Phase 09 &
Tier 2 &
Maintain Tier 2; high volume increases sweat and gear turnover &
Phase 09 (End): sustainment reduces illness/skin breakdown risk; keep inventory stable. \\
Phase 10 &
Tier 2 &
Maintain Tier 2; multi-domain gate weeks increase consequence &
Phase 10 (End): staging prevents forgotten tools that would force substitutions and invalidate evidence. \\
Phase 11 &
Tier 2 &
Maintain Tier 2; recovery sensitivity increases the cost of friction &
Phase 11 (End): keep environment organized and predictable; minimize last-minute changes. \\
Phase 12 &
Tier 3 &
Tier 3 redundancy becomes mandatory by Phase 12 (End) per Stage 6 &
Phase 12 (End): redundancy reduces the chance of compromised test windows due to missing tools or broken storage/sustainment systems. \\
Phase 13 &
Tier 3 &
Maintain Tier 3; consolidation requires low variance and high reliability &
Phase 13 (End): no novelty near final protected windows; sustainment must prevent preventable execution failures. \\
\end{longtblr}

\endgroup

7.10.8 Phase Integration Map Notes:

\begin{itemize}
\item If any phase artifact implies Stage 7.10 Tier 1 is required earlier than Phase 00 (End) per Stage 6, Requires Stage 8 review.
\item Climate control devices are intentionally excluded from tier tables; if a downstream artifact requires them as category items, Requires Stage 8 review.
\end{itemize}

\subsection*{7.10.9 Risks, Contraindications, and Red Flags}
\phantomsection
\addcontentsline{toc}{subsection}{7.10.9 Risks, Contraindications, and Red Flags}

\begin{itemize}
\item Trip/fall hazards:
  \begin{itemize}
  \item clutter, cords, unstable stacking, and blocked walkways increase injury risk; treat as immediate correction items.
  \end{itemize}
\item Tip hazards:
  \begin{itemize}
  \item unanchored tall shelving and top-heavy storage creates fall risk; anchor or replace.
  \end{itemize}
\item Sharp edges and overloaded shelving:
  \begin{itemize}
  \item cuts and collapse risk; reduce load and replace damaged components.
  \end{itemize}
\item Mold/mildew from wet gear storage:
  \begin{itemize}
  \item increases skin irritation and illness risk; drying workflow is mandatory; cross-reference Stage 7.20 for hygiene and Stage 7.14 for wet-weather apparel context.
  \end{itemize}
\item Equipment failure from neglected maintenance:
  \begin{itemize}
  \item loose bolts, frayed straps, broken buckles can cause sudden failure; stop use until repaired/replaced.
  \end{itemize}
\item Lost/forgotten items:
  \begin{itemize}
  \item drives missed sessions and unsafe improvisation; staging checklists and go-bag posture are the controls.
  \end{itemize}
\item Requires Stage 8 review triggers:
  \begin{itemize}
  \item any upstream/downstream artifact demanding climate control devices inside 7.10 tier tables,
  \item any mismatch where a phase demands sustainment capability earlier than Stage 6 Tier timing,
  \item any artifact that treats hygiene consumables as 7.10 tier-table items rather than 7.20.
  \end{itemize}
\end{itemize}

\subsection*{7.10.10 Compliance Checklist — Pass/Fail}
\phantomsection
\addcontentsline{toc}{subsection}{7.10.10 Compliance Checklist — Pass/Fail}

\begin{itemize}
\item \textbf{PASS}/\textbf{FAIL}: Tier 0–Tier 3 tables use required schema, generic item types only, and price bands limited to \$/\$\$/\$\$\$.
\item \textbf{PASS}/\textbf{FAIL}: Tier 0 includes the required household-only examples (single-bin system; designated drying corner; paper checklist/inventory sheet).
\item \textbf{PASS}/\textbf{FAIL}: Climate control devices are not listed in tier tables and are only cross-referenced.
\item \textbf{PASS}/\textbf{FAIL}: Tier 2 includes optional go-bag/vehicle staging concept and does not duplicate 7.20 hygiene consumables.
\item \textbf{PASS}/\textbf{FAIL}: Storage is stable and hazard-controlled (heavy items low; walk paths clear; tip hazards addressed).
\item \textbf{PASS}/\textbf{FAIL}: Drying workflow exists and is used after every session; wet gear never stored sealed while damp.
\item \textbf{PASS}/\textbf{FAIL}: Inventory and maintenance actions can be logged with timestamped entries; missing fields recorded as \textbf{MISSING}/\textbf{UNKNOWN}.
\item \textbf{PASS}/\textbf{FAIL}: Phase Integration Map aligns to Stage 6 timing; mismatches flagged for Stage 8 review (not repaired).
\end{itemize}

\end{document}